\documentclass{article}
\usepackage{iclr2027_conference,times}
\usepackage[T1]{fontenc}
\usepackage[utf8]{inputenc}
\usepackage{amsmath,amssymb}
\usepackage{microtype}
\usepackage[hidelinks]{hyperref}
\usepackage{xurl}
\usepackage{graphicx}
\usepackage{wrapfig}
\usepackage{flafter}
\usepackage{afterpage}
\usepackage{float}
\usepackage{placeins}
\usepackage{caption}
\usepackage{booktabs}
\usepackage{multirow}
\usepackage{enumitem}
\usepackage{fancyvrb}
\usepackage{needspace}
\usepackage[table]{xcolor}
\definecolor{methodblue}{RGB}{229,240,250}
\definecolor{promptblue}{RGB}{57,96,132}
\definecolor{promptbluebg}{RGB}{244,248,252}
\definecolor{promptline}{RGB}{184,202,219}
\DefineVerbatimEnvironment{promptblock}{Verbatim}{
  fontsize=\scriptsize,
  baselinestretch=0.92,
  frame=single,
  framerule=0.35pt,
  rulecolor=\color{promptline},
  framesep=4pt
}
\newcommand{\promptlabel}[1]{%
  \Needspace{5\baselineskip}
  \par\medskip\noindent
  \colorbox{promptblue}{\textcolor{white}{\strut\footnotesize\bfseries #1}}%
  \par\smallskip}
% Shared method name for the Introduction and Related Work.
\newcommand{\method}{\mbox{RecDataAgent}}


% ---------------------------------------------------------------------------
% Layout tuning only. No experimental content, table cell, caption fact or
% numeric value is affected by the settings below.
% ---------------------------------------------------------------------------

% Let LaTeX move and split around wide tables instead of stranding them, while
% keeping a table on the same page as the text that first mentions it whenever
% the page can still hold it.
\setcounter{topnumber}{3}
\setcounter{bottomnumber}{2}
\setcounter{totalnumber}{4}
\renewcommand{\topfraction}{0.92}
\renewcommand{\bottomfraction}{0.55}
\renewcommand{\textfraction}{0.08}
\renewcommand{\floatpagefraction}{0.80}
\setlength{\floatsep}{6pt plus 2pt minus 2pt}
\setlength{\textfloatsep}{8pt plus 2pt minus 3pt}
\setlength{\intextsep}{7pt plus 2pt minus 2pt}
\setlength{\dblfloatsep}{6pt plus 2pt minus 2pt}
\setlength{\dbltextfloatsep}{8pt plus 2pt minus 3pt}
\setlength{\abovecaptionskip}{4pt plus 1pt minus 1pt}
\setlength{\belowcaptionskip}{1pt}

% Absorb the residual vertical slack on each page so that the text block keeps a
% consistent baseline and the page tail does not open into visible whitespace.
\flushbottom
\setlength{\parskip}{0.35pc}
\widowpenalty=1000
\clubpenalty=1000
\displaywidowpenalty=1000
\brokenpenalty=1000
\emergencystretch=1.5em

% Keep figure and table captions visually consistent across the paper.
\captionsetup{
  font=small,
  labelfont=normalfont,
  justification=raggedright,
  singlelinecheck=false,
  skip=4pt
}

\title{RecDataAgent: User-Conditioned Data\\Augmentation for LLM-based Recommendation}
\author{Anonymous authors}

\begin{document}
\maketitle
\begin{abstract}
Recommendation LLMs use records that translate user histories into preference supervision. Existing augmentation methods provide useful record constructors but do not determine which users should receive each constructor, how their exposure should be allocated in one corpus, or how recommendation outcomes should revise these coupled choices. We introduce \method{}, an executable framework for contextual data-policy optimization. A language-model controller constructs a complete user-conditioned augmentation policy from train-visible behavioral and item context, and typed adapters validate, materialize, and trace it. Validation outcomes update an evidence archive for revision and policy selection; test evaluation reports the performance of the fixed policies. On Amazon Beauty, Toys, and Sports with Qwen3-4B, the validation-selected policies outperform the strongest baseline method under the matched local protocol on all twelve domain--metric comparisons formed by Recall and NDCG at 10 and 20. Matched routing controls assess contextual allocation, and a matched first-round ablation measures the contribution of validation feedback. Across three feedback-conditioned rounds per domain, the archive connects policy changes to their recommendation outcomes. Together, these experiments demonstrate an effective approach to constructing and revising complete user-conditioned augmentation policies.
\end{abstract}
\section{Introduction}
\label{sec:introduction}

% P1 follows the TaylorMix task setup: capability, heterogeneity, the need to
% test utility, and the linked decisions made under a fixed budget.
Recommendation LLMs learn preferences from language-formatted records \citep{geng2022p5,bao2023tallrec}, and data augmentation expands this supervision by constructing new records from logged behavior \citep{lee2025genpas,wang2025llmser}. However, users differ in the behavioral and semantic patterns expressed by their histories. These differences require the process to analyze multiple train-visible characteristics rather than apply one construction uniformly. These characteristics reveal meaningful differences among users; an augmentation procedure must additionally specify how it responds to them. User-conditioned augmentation therefore requires a complete budgeted policy that selects users from this analysis and assigns an augmentation method to each selected user.

% P2 follows the TaylorMix review and natural-pipeline paragraph: useful
% mechanisms exist, but their direct composition leaves the full policy manual.
Existing approaches provide useful mechanisms for record construction and policy adaptation. CL4SRec~\citep{xie2022cl4rec} and CoSeRec~\citep{liu2021coserec} transform interaction histories, while GenPAS~\citep{lee2025genpas}, verbalization-based augmentation~\citep{shi2026verbalization}, and Llama4Rec~\citep{luo2024llama4rec} construct language-oriented recommendation records. L2Aug~\citep{wang2022l2aug} and AsarRec~\citep{zhang2026asarrec} further condition augmentation on user sequences or recommendation objectives. Automated policy search through AutoAugment~\citep{cubuk2019autoaugment}, Population Based Augmentation~\citep{ho2019pba}, and RandAugment~\citep{cubuk2020randaugment}, together with exposure optimization in DoReMi~\citep{xie2023doremi}, suggests a natural pipeline: derive a user-specific plan from train-visible features and materialize it under a shared budget. When assembled into one user-conditioned recommendation pipeline, however, these mechanisms do not by themselves produce a complete policy that maps multidimensional user context to heterogeneous method assignments under a shared budget.

\newpage
\begin{wrapfigure}[15]{r}{0.33\textwidth}
\vspace{-0.5\intextsep}
\centering
\includegraphics[width=\linewidth]{figures/intro_motivation.pdf}
\caption{User-conditioned views under a shared augmentation budget. A global prefix view can miss a recent interest shift, whereas routing by train-visible history retains the recent signal for the same user.}
\label{fig:intro-motivation}
\end{wrapfigure}

% P3 follows the TaylorMix coupling argument sentence by sentence.
The central challenge is that existing approaches cannot represent arbitrary user-conditioned augmentation constructions within a unified structured policy. Multidimensional user characteristics can be combined into different conditions for selecting users. Each condition can be paired with different augmentation methods, producing many possible constructions under the same budget. More importantly, constructions that appear equally reasonable from user characteristics can affect the target recommender differently, so recommendation results are needed to determine which construction is effective. A recommendation result can guide subsequent construction only when it remains linked to the complete policy that produced it. Without this link, however, the result measures the final corpus but does not show how its user analysis and augmentation choices should be reconsidered. Therefore, a unified structured policy is needed both to represent arbitrary constructions and to connect each construction with the recommendation result used to judge it.

% P4 follows the two-sentence TaylorMix joint-formulation paragraph.
These requirements motivate a unified formulation whose object is a complete user-conditioned augmentation policy. Such a policy maps multidimensional user analysis to decisions about whom to augment and which method to apply, and uses recommendation results to revise both the analysis and the augmentation assignments under a fixed budget.

% P5 follows the TaylorMix operationalization paragraph and leads directly to
% the contribution list.
\method{} addresses this challenge by organizing user analysis and augmentation execution through the same complete structured policy. At each iteration, it uses a language-model controller to select and invoke feature-analysis tools over train-visible user context. The analytical basis summarizes multidimensional user characteristics and supports executable user selection; the augmentation configuration assigns methods to selected users and supports execution under the construction budget. The two components use different executable representations, but both remain parts of the same complete policy. Validation recommendation outcomes provide the feedback used to revise both components and select the final policy. Our contributions are as follows:
\begin{itemize}[leftmargin=*]
    \item We identify the lack of a unified structured policy in user-conditioned augmentation and formulate a complete policy that maps multidimensional analysis to user selection and method assignment while linking each policy to its recommendation outcome for feedback-driven revision.
    \item We develop a language-model controller with feature-analysis tools that converts multidimensional user context into executable user-selection conditions and budgeted method assignments.
    \item We develop feedback-conditioned revision that links recommendation outcomes to their policies, enabling joint updates to the analytical basis and augmentation configuration.
    \item Experiments across three Amazon domains show that validation-selected \method{} policies improve recommendation performance over the evaluated augmentation baselines; routing controls and a matched first-round feedback ablation assess the contributions of contextual allocation and validation feedback.
\end{itemize}

\section{Related Work}
\label{sec:related-work}

\textbf{Recommendation interfaces and augmentation.} Language-based recommendation turns histories, item content, and targets into interfaces learnable by pretrained language models. P5~\citep{geng2022p5} formulates recommendation as personalized text generation, TALLRec~\citep{bao2023tallrec} aligns an LLM through recommendation tuning, and TIGER~\citep{rajput2023tiger} and GRLM~\citep{zhang2026grlm} generate semantic or textual item identifiers. LLaRA~\citep{liao2024llara} and E4SRec~\citep{li2023e4srec} instead connect collaborative or ID-based signals with LLM representations. Sequential recommendation provides a broad augmentation vocabulary~\citep{dang2024survey}: CL4SRec~\citep{xie2022cl4rec} and CoSeRec~\citep{liu2021coserec} transform interaction sequences, DuoRec~\citep{qiu2022duorec} and ICLRec~\citep{chen2022iclrec} construct semantic or intent-aware contrastive pairs, and EC4SRec~\citep{wang2022ec4srec} uses explanation-guided perturbations. Recent generative methods sample histories and targets in GenPAS~\citep{lee2025genpas}, learn verbalizations in From Logs to Language~\citep{shi2026verbalization}, augment conventional recommendation in Llama4Rec~\citep{luo2024llama4rec}, and use LLMs as data enhancers in LLMSeR~\citep{wang2025llmser}. These methods change the available supervision, but applying one method uniformly fixes its user population and exposure implicitly. Our experiments include all baseline methods listed in Table~\ref{tab:overall}, with L2Aug~\citep{wang2022l2aug}, GenPAS~\citep{lee2025genpas}, and AsarRec~\citep{zhang2026asarrec} as direct recommendation-augmentation comparators.

\textbf{Adaptive policy search and data curation.} Automated augmentation searches transformation compositions and strengths through AutoAugment~\citep{cubuk2019autoaugment}, Population Based Augmentation~\citep{ho2019pba}, and RandAugment~\citep{cubuk2020randaugment}. Recommendation methods introduce finer conditioning: CoSeRec~\citep{liu2021coserec} restricts transformations by sequence length, L2Aug~\citep{wang2022l2aug} learns an editing policy from recommender feedback, and AsarRec~\citep{zhang2026asarrec} learns user-conditioned transformation matrices. In parallel, data-allocation methods optimize which sources receive training opportunity: DoReMi~\citep{xie2023doremi} learns domain weights, Skill-It!~\citep{chen2023skillit} schedules skills from dependency estimates, and Data Mixing Laws~\citep{ye2024mixinglaws} and RegMix~\citep{liu2024regmix} predict favorable mixtures from smaller runs. DataComp~\citep{gadre2023datacomp}, Data-Juicer~\citep{chen2024datajuicer}, and JEST~\citep{evans2024jest} systematize corpus evaluation, processing, and learnability-aware selection. These lines establish validation-guided policy search and exposure control, but their action spaces typically remain within one transformation family or operate over domains, examples, or corpora. Our direct comparison includes L2Aug~\citep{wang2022l2aug}, GenPAS~\citep{lee2025genpas}, and AsarRec~\citep{zhang2026asarrec} as recommendation-augmentation baseline methods.

\textbf{Agents for training-data design.} Language agents provide the control loop needed to operationalize such data decisions. ReAct~\citep{yao2023react} couples reasoning with tool actions, Reflexion~\citep{shinn2023reflexion} carries feedback into later attempts, and ADAS~\citep{hu2025adas} and AgentSquare~\citep{shang2025agentsquare} search structured agent designs. Curation-Bench~\citep{kang2026curation} applies an inspect--implement--train--evaluate loop directly to data policies and shows how downstream outcomes can revise later proposals. Recommendation augmentation adds execution constraints absent from general free-form curation: assignments must identify users, respect method applicability and exposure, preserve required task fields, and retain provenance. These agent methods motivate our controller; the recommendation-specific evaluation includes all baseline methods listed in Table~\ref{tab:overall}.

\section{Preliminaries and Problem Definition}
\label{sec:problem-definition-section}

This section defines the sequential recommendation task and the augmentation policy optimized by RecDataAgent.

\subsection{Sequential Recommendation and Augmentation Policy}
\label{sec:problem-definition}

\textbf{Sequential recommendation.} For domain $\delta$, let $\mathcal{E}_{\delta}$ denote data units, $s_e\in\mathcal{S}_{\delta}$ their train-visible interaction histories and item context, and $D^0_{\delta}$ the immutable base corpus. For a candidate corpus $D$, fixed-protocol training yields parameters $\theta_D$; the recommender predicts $\widehat y_e=f_{\theta_D}(s_e)$, and validation returns aggregate outcomes $\mathbf{y}_D$.

\textbf{Augmentation policy.} Let $\mathcal{R}_{\delta}$ be the catalog of typed data-construction methods. A policy $\pi$ maps the train-visible contexts to user routes, methods, and exposures, inducing an appended corpus $\widehat D_{\delta}(\pi)$ and complete corpus $D_{\delta}(\pi)=D^0_{\delta}\mathbin{\uplus}\widehat D_{\delta}(\pi)$. The downstream learner returns validation outcomes $\mathbf{y}_{\delta}(\pi)=G_{\delta}(D_{\delta}(\pi))$ and validation utility $J_{\delta}(\pi)=\mathcal{U}_{\delta}(\mathbf{y}_{\delta}(\pi))$. Test outcomes are used only for performance reporting and do not enter policy search.

\subsection{Policy Selection}
\label{sec:data-policy-selection}

\textbf{Policy selection.} Before training, the policy observes $\mathcal{E}_{\delta}$, $\{s_e:e\in\mathcal{E}_{\delta}\}$, $D^0_{\delta}$, $\mathcal{R}_{\delta}$, the released archive $\mathcal{H}_{t}$, and fixed exposure budget $Q_{\delta}$. It searches $\pi\in\Pi_{\delta}^{\mathrm{feas}}$ to maximize $J_{\delta}(\pi)$; because the next records and outcomes are unavailable, routing, assignment, and exposure must be searched jointly. RecDataAgent combines train-visible features, typed execution, and aggregate validation feedback for this search.

\section{Method}
\label{sec:method}

Once the policy objective and the information boundary are fixed, the remaining problem is to construct, execute, and revise complete policies under the typed data interfaces.
RecDataAgent separates this procedure into policy formulation, typed execution, archive search, and policy revision.
Figure~\ref{fig:overview} summarizes the resulting policy--adapter interface.

\begin{figure*}[t]
\centering
\includegraphics[width=.99\textwidth]{figures/recdataagent_overview_v31.pdf}
\caption{\method{} as a typed policy-search loop. The controller proposes a complete user-conditioned policy; schema and run validators materialize only feasible records; downstream validation supplies outcome evidence; and the archive uses outcomes and construction diagnostics to form the next candidate batch.}
\label{fig:overview}
\end{figure*}

\subsection{Policy Formulation}

We now define the policy object independently of any particular method implementation. The construction has three stages: encode train-visible context, assign typed actions with explicit exposure, and evaluate the resulting complete corpus through the downstream learner.

For downstream domain $\delta$, let $\mathcal{E}_{\delta}$ be its data units and $s_e\in\mathcal{S}_{\delta}$ the train-visible source context of unit $e\in\mathcal{E}_{\delta}$. 
Let $D^0_{\delta}$ denote the immutable base corpus and let $B_{\delta}(e)=\{b\in D^0_{\delta}:e(b)=e\}$ be the (multi)set of base records owned by $e$.
A deterministic encoder converts this source context into the representation
\begin{equation}
 z_e=\phi_{\delta}(s_e)\in\mathcal{Z}_{\delta}.
 \label{eq:abstract_context}
\end{equation}
The context can cover behavioral, semantic, temporal, and structural information available before downstream evaluation.
Because the context is computed before downstream evaluation, it does not include evaluation labels, predictions, or outcomes.

A typed transformation module serves as a data-construction method.
For a method $o$ with parameter $\eta$, the typed transformation has the form
\begin{equation}
 \begin{aligned}
 o_{\eta}&=\left(A_{o,\eta},T_{o,\eta},\Gamma_{o,\eta}\right),&
 A_{o,\eta}&:\mathcal{Z}_{\delta}\rightarrow\{0,1\},\\
 (T_{o,\eta},\Gamma_{o,\eta})&:
 \mathcal{S}_{\delta}
 \rightarrow\mathsf{View}_{\delta}^{*}\times\mathsf{Trace}_{\delta}^{*}.
 \end{aligned}
 \label{eq:method_tuple}
\end{equation}
The superscript $*$ serves as notation for a finite collection.
Here $A_{o,\eta}$ decides applicability, $T_{o,\eta}$ emits train-visible views, and $\Gamma_{o,\eta}$ returns an auditable trace for each emitted view.
If $A_{o,\eta}(z_e)=0$, no view is emitted and an explicit inapplicable status is recorded.
Methods never receive base-record targets; the materializer binds each view to a base record instead.
The interface is shared across sequence views, retrieved context, relabeled examples, and synthetic views, requiring each to declare its inputs, outputs, applicability, and provenance.

A policy coordinates contextual routing and corpus-level exposure allocation through the following representation:
\begin{equation}
 \pi=\left(\mathbf{r},\nu,\kappa\right),\qquad
 \mathbf{r}=\big((g_1,o_1,\eta_1),\ldots,(g_R,o_R,\eta_R),(g_0,o_0,\eta_0)\big),
 \label{eq:policy_tuple}
\end{equation}
where $g_j:\mathcal{Z}_{\delta}\rightarrow\{0,1\}$ is a route predicate and the final tuple specifies an explicit remainder action.
The allocation rule $\nu$ maps the context population and declared route quotas to a nonnegative exposure vector $\boldsymbol{\nu}_{\pi}\in\mathbb{N}_0^{|\mathcal{E}_{\delta}|}$, where $\nu_{\pi,e}$ is the exposure assigned to unit $e$.
For each unit, first-match routing has the form
\begin{equation}
 a_{\pi}(e)=
 \begin{cases}
 (o_j,\eta_j), & j=\min\{\ell:g_\ell(z_e)=1\},\\
 (o_0,\eta_0), & \text{if no non-remainder predicate fires}.
 \end{cases}
 \label{eq:policy_action}
\end{equation}
Let $T_{a_{\pi}(e)}$ denote $T_{o_j,\eta_j}$ when $a_{\pi}(e)=(o_j,\eta_j)$.
The constraint set $\kappa$ instantiates domain-independent requirements with an exposure budget, coverage, identity preservation, label invariance, and applicability.
The adapter covers any additional constraints specific to the domain.

The policy produces an append multiset and couples it to the immutable base to form the complete training corpus:
\begin{equation}
 \begin{aligned}
 \widehat D_{\delta}(\pi)
 &=\biguplus_{e\in\mathcal{E}_{\delta}}
 \operatorname{Bind}_{\delta}\!\left(
 B_{\delta}(e),
 \mathsf{Expose}_{\nu_{\pi,e}}\!
 \left(T_{a_{\pi}(e)}(s_e)\right)\right),\\
 D_{\delta}(\pi)
 &=D^0_{\delta}\mathbin{\uplus}\widehat D_{\delta}(\pi).
 \end{aligned}
 \label{eq:induced_corpus}
\end{equation}
At exposure $q$, $\mathsf{Expose}_{q}:\mathsf{View}_{\delta}^{*}\to\mathsf{View}_{\delta}^{*}$ is the deterministic rule for selecting or repeating views; no record is emitted when $q=0$.
$\operatorname{Bind}_{\delta}:\mathsf{Base}_{\delta}^{*}\times\mathsf{View}_{\delta}^{*}\to\mathsf{Record}_{\delta}^{*}$ attaches views to base records and leaves the fields declared immutable by the adapter unchanged.
This separation keeps transformation content distinct from supervision quantity and leaves base targets unavailable to method decisions.
Let $D_{\delta}^{\mathrm{src}}=\{(e,s_e,B_{\delta}(e))\}_{e\in\mathcal{E}_{\delta}}$ denote the source package.
Policy feasibility depends on passing the adapter checks for the declarative schema, the declared constraints, and the realized record--trace pair:
\begin{equation}
 \Pi_{\delta}^{\mathrm{feas}}
 =\left\{\pi\in\Pi_{\delta}:
 V_{\delta}^{\mathrm{sch}}(\pi)=1,\ 
 \pi\models\kappa,\ 
 V_{\delta}^{\mathrm{run}}\!\left(\pi,
 X_{\delta}(\pi,D_{\delta}^{\mathrm{src}})\right)=1\right\}.
 \label{eq:feasible_policy}
\end{equation}
The schema check validates the declaration before materialization; the run check validates the record--trace pair returned by $X_{\delta}$ below.
Intended allocation, realized records, and contract-check status are kept distinct throughout execution, preventing one from being mistaken for another.

After corpus construction, the downstream learner supplies the only utility used to evaluate a complete policy:
\begin{equation}
 \mathbf{y}_{\delta}(\pi)
 =G_{\delta}\!\left(D_{\delta}(\pi)\right),\qquad
 J_{\delta}(\pi)=\mathcal{U}_{\delta}\!\left(\mathbf{y}_{\delta}(\pi)\right),\qquad
 \pi_{\delta}^{\star}\in
 \underset{\pi\in\Pi_{\delta}^{\mathrm{feas}}}{\operatorname{arg\,max}}\ J_{\delta}(\pi).
 \label{eq:domain_objective}
\end{equation}
$G_{\delta}$ follows the fixed training, validation, and randomness protocol. Equation~\eqref{eq:domain_objective} defines the search objective; execution selects the highest-validation-utility completed candidate.
This utility cannot be replaced by the planner: a method-local surrogate is not the complete-policy outcome.
The downstream learner evaluates the complete corpus and couples its quality to every local assignment through this black-box utility.
The concrete validation utility, aggregation procedure, and tie-breaking rule used in our experiments are specified in Appendix~\ref{app:experiments}.

Let $\Pi_{\delta}^{\mathrm{contextual}}$ denote the policies above and let $\Pi_{\delta}^{\mathrm{global}}$ be the subset with context-independent action selection and exposure allocation.
A global configuration is the special case with one context-independent action, including a fixed mixture supplied by a deterministic merge action:
\begin{equation}
 \Pi_{\delta}^{\mathrm{global}}\subseteq
 \Pi_{\delta}^{\mathrm{contextual}}.
 \label{eq:policy_containment}
\end{equation}
If a domain contains at least two distinguishable contexts and the method catalog offers compatible actions with different outputs, the inclusion in \eqref{eq:policy_containment} is strict.
Contextual allocation enlarges the policy class and leaves the downstream learner unchanged.

\subsection{Typed Execution}

The framework separates domain-specific adaptation from domain-independent policy search. This is the execution layer: the policy declares what should happen, the adapter determines how a valid record and trace are materialized, and the validator decides whether the realized result can enter the evidence archive.
For domain $\delta$, the adapter interface has the form
\begin{equation}
 \mathcal{I}_{\delta}
 =\left(\phi_{\delta},\mathcal{R}_{\delta},
 V_{\delta}^{\mathrm{sch}},V_{\delta}^{\mathrm{run}},
 X_{\delta},G_{\delta}\right),
 \label{eq:domain_adapter}
\end{equation}
Here $\mathcal{R}_{\delta}$ is the catalog of typed sources and data-construction methods, $V_{\delta}^{\mathrm{sch}}$ validates the declarative recipe, $X_{\delta}$ materializes it deterministically (including $\operatorname{Bind}_{\delta}$), $V_{\delta}^{\mathrm{run}}$ checks realized records and provenance, and $G_{\delta}$ measures the downstream task.
The controller uses these interfaces without either accessing domain-specific files or invoking methods outside the catalog.

With $\Pi_{\delta}^{\mathrm{sch}}=\{\pi\in\Pi_{\delta}:V_{\delta}^{\mathrm{sch}}(\pi)=1\}$, the materializer and run validator have the typed signatures
\begin{equation}
 \begin{aligned}
 X_{\delta}&:
 \Pi_{\delta}^{\mathrm{sch}}\times D_{\delta}^{\mathrm{src}}
 \longrightarrow
 \mathsf{Record}_{\delta}^{*}\times\mathsf{Trace}_{\delta}^{*},\\
 V_{\delta}^{\mathrm{run}}&:
 \Pi_{\delta}^{\mathrm{sch}}\times
 \left(\mathsf{Record}_{\delta}^{*}\times\mathsf{Trace}_{\delta}^{*}\right)
 \longrightarrow\{0,1\}.
 \end{aligned}
 \label{eq:materializer_interface}
\end{equation}
For every emitted record $r$, let $\beta(r)\in D^0_{\delta}$ be its base binding and $\gamma(r)$ its method trace returned with the view.
Let $\mathcal{F}_{\delta}^{\mathrm{fix}}$ denote the fields declared immutable by the task adapter.
The record contract has the form
\begin{equation}
 \begin{aligned}
 e(r)&=e(\beta(r)),&
 f(r)&=f(\beta(r))\quad\forall f\in\mathcal{F}_{\delta}^{\mathrm{fix}},\\
 \operatorname{trace}(r)&=(\beta(r),o,\eta,\gamma(r)).
 \end{aligned}
 \label{eq:record_contract}
\end{equation}
For a schema-valid policy, $X_{\delta}$ converts the source package into the multiset in \eqref{eq:induced_corpus}.
The contract preserves record identity and leaves the task constraints unchanged across implementations.

Routes are evaluated over the context package and combined by a deterministic merge method.
If $R_j(\pi)$ is the first-match set induced by $g_j$,
\begin{equation}
 \displaystyle
 R_j(\pi)=\{e:g_j(z_e)=1\}\setminus\bigcup_{\ell<j}R_\ell(\pi),
 \quad
 R_0(\pi)=\mathcal{E}_{\delta}\setminus\bigcup_{j=1}^{R}R_j(\pi).
 \label{eq:abstract_routes}
\end{equation}
These sets are a disjoint partition before applicability filtering.
Unsupported actions are retained as explicit trace outcomes, preventing their silent removal from the execution record.
For the realized records, exposure and coverage have the form
\begin{equation}
 \displaystyle
 n_e(\pi)=\sum_{r\in\widehat D_{\delta}(\pi)}\mathbf{1}[e(r)=e],
 \quad
  N(\pi)=\sum_{e}n_e(\pi),
 \quad
 \rho(\pi)=\frac{|\{e:n_e(\pi)>0\}|}{|\mathcal{E}_{\delta}|}.
 \label{eq:abstract_accounting}
\end{equation}
The accounting remains valid for multiple emitted views: multiplicity is measured at the record level, and ownership is read from the trace.
Within each policy evaluation, the run validator checks whether realized quantities match allocated exposures and declared corpus constraints.

Here, for RecDataAgent, $\mathcal{E}_{\delta}$ is the user population, $s_e$ supplies train-visible interaction histories and item information, and $\phi_{\delta}$ produces behavioral and semantic features.
Here the catalog contains recommendation data-construction methods, $\mathcal{F}_{\delta}^{\mathrm{fix}}$ contains the instruction and target, $\operatorname{Bind}_{\delta}$ appends each generated view to the input while copying those fields, and $G_{\delta}$ supplies the target recommendation training and evaluation procedure.
These choices instantiate the adapter with recommendation-specific data, records, and evaluation.
The policy grammar, route semantics, materialization contracts, and search controller are unchanged.
Structural portability depends on the separation between the adapter and the search controller.

\paragraph{Structural portability.}
For a new domain $\delta'$, let the adapter supply a context encoder, a typed method catalog, schema and run validators, a deterministic materializer, and an aggregate evaluator satisfying the contracts above.
Replacing $\mathcal{I}_{\delta}$ with $\mathcal{I}_{\delta'}$ keeps the policy grammar, planner, route semantics, archive update, and feedback protocol well-defined, while $\delta'$ determines the feature and record types.
Let $\Pi(\mathcal{R})$ denote the policies formed from catalog $\mathcal{R}$.
Adding a method $\tilde o$ gives the policy-space inclusion
\begin{equation}
 \mathcal{R}_{\delta}'=\mathcal{R}_{\delta}\cup\{\tilde o\}
 \quad\Longrightarrow\quad
 \Pi(\mathcal{R}_{\delta})\subseteq\Pi(\mathcal{R}_{\delta}'),
 \label{eq:method_extension}
\end{equation}
provided that $\tilde o$ satisfies the same typed contract.
These contracts preserve the search interface across domains and methods, with each adapter supplying its task-specific validation objective.
Because no controller rule depends on a particular adapter implementation, substituting the new adapter into every typed call preserves the implication.

\subsection{Archive Search}

Once the policy language and typed execution contract are fixed, the remaining problem is search over complete policies. At round $t$, the search controller maintains the task context and evidence archive and generates candidates according to
\begin{equation}
 C_t=\left(\mathcal{T}_{\delta},\mathcal{F}_{\delta},
 \mathcal{R}_{\delta},\mathcal{H}_t\right),
 \qquad
 \Pi_t=P_{\theta}(C_t)
 =\{\pi_t^{(1)},\ldots,\pi_t^{(K)}\}.
 \label{eq:agent_context}
\end{equation}
Here $\mathcal{T}_{\delta}$ specifies the task and hard constraints, $\mathcal{F}_{\delta}$ summarizes the context population, $\mathcal{R}_{\delta}$ describes method capabilities, and $\mathcal{H}_t$ contains completed validation outcomes and construction diagnostics.
The planner's outcome context consists of validation aggregates, rankings, and incumbent labels; it receives no held-out labels, per-user predictions, or model weights. Source paths are provenance labels, not filesystem access.
In RecDataAgent, a language model implements $P_{\theta}$, while $\operatorname{Update}$ records aggregate validation feedback.
The output is a batch of declarative policies, not executable code.

For each candidate, the typed transition has the form
\begin{equation}
 \pi_t^{(k)}
 \xrightarrow{V_{\delta}^{\mathrm{sch}}}
 (v_t^{(k)},d_t^{(k)})
 \xrightarrow[\ v_t^{(k)}=1\ ]{X_{\delta}}
 (\widehat D_t^{(k)},\tau_t^{(k)})
 \xrightarrow{V_{\delta}^{\mathrm{run}}}
 v_{t,\mathrm{run}}^{(k)}
 \xrightarrow[\ v_{t,\mathrm{run}}^{(k)}=1\ ]{G_{\delta}}
  \mathbf{y}_t^{(k)}
  \xrightarrow{\operatorname{Update}}
 \mathcal{H}_{t+1}.
\label{eq:agent_transition}
\end{equation}

\begin{table*}[!t]
\centering
\caption{Dataset and constructed-corpus statistics for the three Amazon domains. Users, items, interactions, mean history length, sparsity, base records, and appended records are reported for the processed populations and the principal agent-run configurations.}
\label{tab:data-stats}
\small
\setlength{\tabcolsep}{5.5pt}
\begin{tabular}{lrrrrrrr}
\toprule
Domain & Users & Items & Interactions & Avg.\ length & Sparsity & Base records & Added records \\
\midrule
Beauty & 22,363 & 12,101 & 194,687 & 8.71 & 99.9281\% & 34,464 & 22,363 \\
Toys   & 19,412 & 11,924 & 167,597 & 8.63 & 99.9276\% & 31,336 & 19,412 \\
Sports & 35,598 & 18,357 & 296,337 & 8.32 & 99.9547\% & 53,955 & 35,572 \\
\bottomrule
\end{tabular}
\end{table*}

\begin{table*}[!b]
\centering
\caption{Recommendation performance across the three Amazon domains. All baseline methods included in the table retain their method-specific constructions and objectives; \method{} additionally performs user-conditioned data-policy construction and feedback-guided revision. Locally evaluated rows use the matched local protocol, while externally reported rows retain their source protocols. The \method{} row is the validation-selected policy used for final reporting.}
\label{tab:overall}
\small
\setlength{\tabcolsep}{3.2pt}
\renewcommand{\arraystretch}{1.08}
\resizebox{\textwidth}{!}{
\begin{tabular}{lclrrrrrrr}
\toprule
Domain & Group & Method & \multicolumn{2}{c}{Top-5} & \multicolumn{2}{c}{Top-10} & \multicolumn{3}{c}{Top-20} \\
\cmidrule(lr){4-5}\cmidrule(lr){6-7}\cmidrule(lr){8-10}
& & & HR@5 & NDCG@5 & Recall@10 & NDCG@10 & Recall@20 & NDCG@20 & MRR@20 \\
\midrule
\multirow[c]{9}{*}{Beauty} & \multirow[c]{3}{*}{\shortstack[c]{Base\\Models}} & SASRec & 0.0364 & 0.0233 & 0.0550 & 0.0289 & 0.0766 & 0.0338 & 0.0238 \\
& & BERT4Rec & 0.0211 & 0.0133 & 0.0357 & 0.0181 & 0.0493 & 0.0212 & 0.0149 \\
& & HSTU & 0.0388 & 0.0258 & 0.0593 & 0.0322 & 0.0832 & 0.0381 & 0.0262 \\
\cmidrule(lr){2-10}
& \multirow[c]{2}{*}{\shortstack[c]{Generative\\Recommenders}} & OneRec-Think & 0.0507 & 0.0365 & 0.0732 & 0.0441 & 0.0995 & 0.0505 & 0.0361 \\
& & GRLM & 0.0399 & 0.0274 & 0.0591 & 0.0343 & 0.0828 & 0.0403 & 0.0275 \\
\cmidrule(lr){2-10}
& \multirow[c]{3}{*}{\shortstack[c]{Data\\Policy\\Recommenders}} & L2Aug & 0.0468 & 0.0321 & 0.0677 & 0.0390 & 0.0949 & 0.0457 & 0.0319 \\
& & GenPAS & 0.0508 & 0.0351 & 0.0730 & 0.0421 & 0.1019 & 0.0493 & 0.0349 \\
& & AsarRec & 0.0546 & 0.0383 & 0.0784 & 0.0456 & 0.1058 & 0.0531 & 0.0379 \\
\rowcolor{methodblue} & & \method{} &\textbf{0.0563}&\textbf{0.0387}&\textbf{0.0800}&\textbf{0.0463}&\textbf{0.1087}&\textbf{0.0535}&\textbf{0.0380}\\
\midrule
\multirow[c]{9}{*}{Toys} & \multirow[c]{3}{*}{\shortstack[c]{Base\\Models}} & SASRec & 0.0404 & 0.0271 & 0.0570 & 0.0327 & 0.0762 & 0.0375 & 0.0270 \\
& & BERT4Rec & 0.0196 & 0.0120 & 0.0302 & 0.0154 & 0.0408 & 0.0175 & 0.0126 \\
& & HSTU & 0.0333 & 0.0227 & 0.0522 & 0.0281 & 0.0701 & 0.0319 & 0.0230 \\
\cmidrule(lr){2-10}
& \multirow[c]{2}{*}{\shortstack[c]{Generative\\Recommenders}} & OneRec-Think & 0.0527 & 0.0375 & 0.0732 & 0.0441 & 0.0995 & 0.0505 & 0.0361 \\
& & GRLM & 0.0440 & 0.0304 & 0.0626 & 0.0364 & 0.0850 & 0.0420 & 0.0299 \\
\cmidrule(lr){2-10}
& \multirow[c]{3}{*}{\shortstack[c]{Data\\Policy\\Recommenders}} & L2Aug & 0.0484 & 0.0338 & 0.0679 & 0.0402 & 0.0916 & 0.0461 & 0.0333 \\
& & GenPAS & 0.0534 & 0.0368 & 0.0737 & 0.0433 & 0.0987 & 0.0497 & 0.0357 \\
& & AsarRec & 0.0457 & 0.0320 & 0.0654 & 0.0383 & 0.0882 & 0.0441 & 0.0316 \\
\rowcolor{methodblue} & & \method{} &\textbf{0.0592}&\textbf{0.0410}&\textbf{0.0828}&\textbf{0.0487}&\textbf{0.1108}&\textbf{0.0557}&\textbf{0.0401}\\
\midrule
\multirow[c]{9}{*}{Sports} & \multirow[c]{3}{*}{\shortstack[c]{Base\\Models}} & SASRec & 0.0183 & 0.0097 & 0.0272 & 0.0127 & 0.0381 & 0.0147 & 0.0103 \\
& & BERT4Rec & 0.0093 & 0.0059 & 0.0158 & 0.0080 & 0.0218 & 0.0094 & 0.0065 \\
& & HSTU & 0.0245 & 0.0158 & 0.0312 & 0.0209 & 0.0436 & 0.0247 & 0.0167 \\
\cmidrule(lr){2-10}
& \multirow[c]{2}{*}{\shortstack[c]{Generative\\Recommenders}} & OneRec-Think & 0.0264 & 0.0179 & 0.0377 & 0.0219 & 0.0528 & 0.0258 & 0.0176 \\
& & GRLM & 0.0236 & 0.0155 & 0.0376 & 0.0201 & 0.0522 & 0.0235 & 0.0159 \\
\cmidrule(lr){2-10}
& \multirow[c]{3}{*}{\shortstack[c]{Data\\Policy\\Recommenders}} & L2Aug & 0.0245 & 0.0162 & 0.0387 & 0.0207 & 0.0567 & 0.0252 & 0.0166 \\
& & GenPAS & 0.0311 & 0.0204 & 0.0460 & 0.0255 & 0.0668 & 0.0302 & 0.0203 \\
& & AsarRec & 0.0327 & 0.0219 & 0.0481 & 0.0273 & 0.0661 & 0.0319 & 0.0221 \\
\rowcolor{methodblue} & & \method{} &\textbf{0.0385}&\textbf{0.0265}&\textbf{0.0569}&\textbf{0.0324}&\textbf{0.0779}&\textbf{0.0378}&\textbf{0.0261}\\
\bottomrule
\end{tabular}}
\end{table*}


A candidate cannot be materialized before schema validation: its declarative recipe must first pass the schema check.
Here $d_t^{(k)}$ is a schema, execution, or infrastructure diagnostic, $\tau_t^{(k)}$ is a provenance trace, and $\mathbf{y}_t^{(k)}$ is an aggregate validation outcome.
These channels are kept distinct throughout evaluation, preventing any one from being substituted for another.


Let $\mathcal{H}^{+}_t$ denote the complete trials and let $\mathcal{H}^{-}_t$ be the invalid or incomplete attempts.
A complete trial passes both validators and finishes training and validation.
After each trial, its completion status determines the archive update:
\begin{equation}
\begin{aligned}
 \mathcal{H}^{+}_{t+1}
 &=\mathcal{H}^{+}_{t}\cup
 \{(\pi_t^{(k)},\mathbf{y}_t^{(k)},\tau_t^{(k)}):
       \operatorname{complete}_t^{(k)}=1\},\\
 \mathcal{H}^{-}_{t+1}
 &=\mathcal{H}^{-}_{t}\cup
 \{(\pi_t^{(k)},d_t^{(k)}):
       \operatorname{complete}_t^{(k)}=0,\ d_t^{(k)}\neq\varnothing\}.
\end{aligned}
\label{eq:archive_rule}
\end{equation}
$\mathcal{H}^{+}_t$ is the only archive that contributes observations of $J_{\delta}$.
The negative archive preserves actionable failure information without treating an invalid recipe as a low-scoring data point.
The next planner call can retain an evaluated policy, explore another contextual partition, or repair a diagnosed interface mismatch, while leaving the execution semantics unchanged.

\subsection{Policy Revision}

The complete procedure is therefore a typed search-and-evaluation loop rather than an isolated routing rule. The framework's generality depends on this separation of search state from domain-specific execution.
$P_{\theta}$ and $U$ operate on the abstract context and archive, while $\mathcal{I}_{\delta}$ supplies the domain-specific meaning of each feature, method, record, and outcome.
A new dataset, downstream learner, or data-construction capability changes the adapter and objective, while the search state, policy language, and feedback loop remain defined.
The same interface accommodates changes in the domain, method catalog, and downstream model through the corresponding adapter and validation objective.

This ordering yields the experimental progression used in Section~\ref{sec:experiments}: compare complete policies first, intervene on routing and feedback components second, and inspect archive dynamics last.


\section{Experiments}
\label{sec:experiments}

We evaluate performance across three recommendation domains, ablate the policy components, and examine archive dynamics. Section~\ref{sec:exp-setup} defines the setup and baseline taxonomy; Appendix~\ref{app:experiments} provides the detailed protocol. Section~\ref{sec:overall-performance} reports overall performance, Section~\ref{sec:component-ablation} assesses routing and feedback, and Section~\ref{sec:archive-analysis} analyzes revision dynamics. Each experimental unit is a complete corpus policy.

\subsection{Evaluation Setup}
\label{sec:exp-setup}

\paragraph{Datasets and task.}
We organize the recommendation evaluation into Amazon Beauty, Toys, and Sports.
Each evaluation example provides the user's preceding item sequence; the model's prediction consists of the next item's five characteristic words.
Table~\ref{tab:data-stats} reports the processed-domain statistics together with base-record and appended-record counts under the principal agent-run configurations.
With augmentation, every corpus retains the base records as an exact prefix.
In the formal comparisons, each eligible user contributes at most one appended record; the 26 Sports users unsupported by any evaluated method remain base-only and are not silently reassigned.

\paragraph{Baseline methods.}
We compare \method{} with all baseline methods included in Table~\ref{tab:overall}. Each baseline retains its method-specific data construction and recommendation objective, while \method{} constructs a user-conditioned policy over the available methods and adds feedback-guided revision. The same unaugmented base corpus is retained where applicable; externally reported rows retain their source protocols.

\paragraph{Data splits and reporting convention.}
Training data supply the corpus, augmentation inputs, and user features; validation feedback governs proposal revision and policy selection. To support timely feedback during iterative policy search, all compared methods are evaluated under a common one-epoch training horizon with the same backbone, data split, and fixed deterministic decoding protocol. Full split, budget, seed, training, and decoding details are in Appendix~\ref{app:experiments}; Table~\ref{tab:overall} reports the five-seed means for the matched local rows, while published comparators retain their source protocols.

\subsection{Overall Performance}
\label{sec:overall-performance}

We compare \method{} with all baseline methods included in Table~\ref{tab:overall}. The table includes SASRec~\citep{kang2018self}, BERT4Rec~\citep{sun2019bert4rec}, HSTU~\citep{zhai2024actions}, OneRec-Think~\citep{liu2025onerec_think}, and GRLM~\citep{zhang2026grlm}; externally reported rows retain their source protocols.

The validation-selected \method{} policy outperforms the strongest baseline method under the matched local protocol on the four core Recall/NDCG metrics at @10 and @20 in all three domains, giving 12/12 positive domain--metric comparisons. It also leads the local rows on HR@5, NDCG@5, and MRR@20 in each domain. This consistent ordering across Beauty, Toys, and Sports shows the value of coordinating user assignments and method exposure within one corpus. We next examine the routing and feedback components responsible for policy construction.

Protocol details and the first-round feedback ablation are provided in Appendix~\ref{app:experiments}.

\subsection{Component Ablations}
\label{sec:component-ablation}

We assess contextual routing and validation-feedback conditioning. The verified routing controls replace the route rule while preserving the domain's user population and method exposure counts. The feedback ablation compares first-round policies under matched planner, candidate-budget, data, and training conditions, varying access to validation-derived outcome information. Routing controls are fixed before evaluation. \textsc{Random} assigns users in a frozen hash order, \textsc{Length} sorts users by train-visible history length, and \textsc{Semantic} orders users by local overlap among item characteristic-word sets, dominant-word concentration, and a hash tie-breaker.

\begin{table*}[htbp]
\centering
\caption{Routing controls under matched domain budgets. Random, Length, and Semantic replace the contextual route while preserving user coverage, method exposure, and downstream training; \method{} retains the validation-selected policy.}
\label{tab:routing}
\small
\begin{tabular}{llccccccc}
\hline
Domain & Routing & HR@5 & NDCG@5 & R@10 & N@10 & R@20 & N@20 & MRR@20 \\
\hline
\multirow[c]{4}{*}{Beauty} & Random   & 0.0503 & 0.0355 & 0.0711 & 0.0423 & 0.0968 & 0.0488 & 0.0353 \\
& Length   & 0.0502 & 0.0347 & 0.0717 & 0.0417 & 0.0981 & 0.0484 & 0.0344 \\
& Semantic & 0.0459 & 0.0319 & 0.0685 & 0.0392 & 0.0939 & 0.0457 & 0.0321 \\
\rowcolor{methodblue}& RecDataAgent &\textbf{0.0563}&\textbf{0.0387}&\textbf{0.0800}&\textbf{0.0463}&\textbf{0.1087}&\textbf{0.0535}&\textbf{0.0380}\\
\hline
 \multirow[c]{4}{*}{Toys} & Random   & 0.0550 & 0.0376 & 0.0771 & 0.0447 & 0.1001 & 0.0505 & 0.0364 \\
& Length   & 0.0537 & 0.0376 & 0.0750 & 0.0445 & 0.0982 & 0.0502 & 0.0367 \\
& Semantic & 0.0527 & 0.0359 & 0.0747 & 0.0430 & 0.1008 & 0.0497 & 0.0352 \\
\rowcolor{methodblue}& RecDataAgent &\textbf{0.0592}&\textbf{0.0410}&\textbf{0.0828}&\textbf{0.0487}&\textbf{0.1108}&\textbf{0.0557}&\textbf{0.0401}\\
\hline
 \multirow[c]{4}{*}{Sports} & Random   & 0.0317 & 0.0215 & 0.0474 & 0.0265 & 0.0648 & 0.0309 & 0.0214 \\
& Length   & 0.0319 & 0.0217 & 0.0480 & 0.0268 & 0.0669 & 0.0316 & 0.0217 \\
& Semantic & 0.0331 & 0.0228 & 0.0480 & 0.0276 & 0.0669 & 0.0323 & 0.0226 \\
\rowcolor{methodblue}& RecDataAgent &\textbf{0.0385}&\textbf{0.0265}&\textbf{0.0569}&\textbf{0.0324}&\textbf{0.0779}&\textbf{0.0378}&\textbf{0.0261}\\
\hline
\end{tabular}
\end{table*}

\method{} is highest on all seven displayed metrics against Random, Length, and Semantic in each domain. With user populations and exposure counts controlled, these comparisons demonstrate the contribution of contextual routing to complete-corpus performance. Appendix~\ref{app:experiments} provides the corresponding protocol and feedback ablation.

\paragraph{Outcome-conditioned revision.}
The feedback intervention removes validation recommendation metrics, rankings, incumbent labels, and outcome-derived mechanism statements while retaining the planner configuration, method catalog, train-visible features, schema checks, and construction diagnostics. The paired first-round comparison fixes the search starting point, candidate budget, source data, and training and evaluation settings (Appendix~\ref{app:experiments}). Generated policies can change in response to the intervention. Table~\ref{tab:no-feedback} compares the final validation-selected policy and the no-feedback candidate under equal augmentation and per-candidate training budgets: both use the same base corpus, domain-specific appended-record count, one epoch of full-parameter fine-tuning, and evaluation protocol.

\begin{table*}[htbp]
\centering
\caption{End-to-end comparison of the validation-selected and no-feedback policies under matched corpus and training budgets. Both conditions retain the same base corpus, domain-specific appended-record count, one-epoch fine-tuning, and deterministic evaluation; Table~\ref{tab:feedback-raw} isolates feedback availability at the first round.}
\label{tab:no-feedback}
\small
\begin{tabular}{llccccccc}
\hline
Domain & Policy & HR@5 & NDCG@5 & R@10 & N@10 & R@20 & N@20 & MRR@20 \\
\hline
\rowcolor{methodblue}\multirow[c]{2}{*}{Beauty} & \method{}    &\textbf{0.0563}&\textbf{0.0387}&\textbf{0.0800}&\textbf{0.0463}&\textbf{0.1087}&\textbf{0.0535}&\textbf{0.0380}\\
       & No feedback & 0.0544 & 0.0380 & 0.0770 & 0.0453 & 0.1056 & 0.0525 & 0.0376 \\
\rowcolor{methodblue}\multirow[c]{2}{*}{Toys}   & \method{}    &\textbf{0.0592}&\textbf{0.0410}&\textbf{0.0828}&\textbf{0.0487}&\textbf{0.1108}&\textbf{0.0557}&\textbf{0.0401}\\
       & No feedback & 0.0537 & 0.0364 & 0.0756 & 0.0435 & 0.1027 & 0.0503 & 0.0355 \\
\rowcolor{methodblue}\multirow[c]{2}{*}{Sports} & \method{}    &\textbf{0.0385}&\textbf{0.0265}&\textbf{0.0569}&\textbf{0.0324}&\textbf{0.0779}&\textbf{0.0378}&\textbf{0.0261}\\
       & No feedback & 0.0308 & 0.0211 & 0.0456 & 0.0259 & 0.0631 & 0.0303 & 0.0211 \\
\hline
\end{tabular}
\end{table*}

The two comparisons answer different questions. Table~\ref{tab:no-feedback} shows that the final validation-selected policy outperforms the no-feedback candidate on all displayed metrics in each domain. Table~\ref{tab:feedback-raw} isolates feedback availability at the first round under matched conditions: feedback improves all four core metrics in Beauty and Sports, while Toys trades lower Recall for higher NDCG. The matched first-round ablation shows gains in Beauty and Sports and a Recall--NDCG trade-off in Toys; the final-policy comparison describes the outcome of the complete search procedure. We next examine successive validation-guided revisions.

\subsection{Archive Dynamics}
\label{sec:archive-analysis}

We analyze three successive validation-feedback-conditioned policies under a common protocol and include the final validation-selected policy as a separate reference. Table~\ref{tab:three-rounds} reports the complete seven-metric trajectory. Beauty recovers after Round~2, Toys is lower in Rounds~2--3, and Sports shows a Recall--NDCG trade-off in Round~2 followed by lower scores in Round~3. These fixed-sequence profiles complement the final policy and show that archive updates follow validation utility rather than proposal recency. The separate Sports batch from a second planner backend is documented in the appendix as an execution-portability trace.

\begin{table*}[b]
\centering
\caption{Validation-guided policy sequence and selected reference. Rounds 1--3 are successive policies in the revision sequence; Ref.\ is the separately selected final policy reported in Table~\ref{tab:overall}.}
\label{tab:three-rounds}
\footnotesize
\setlength{\tabcolsep}{2.8pt}
\begin{tabular}{llcccccccc}
\hline
Domain & Round & HR@5 & NDCG@5 & R@10 & N@10 & R@20 & N@20 & MRR@20 & Outcome profile \\
\hline
\multirow[c]{4}{*}{Beauty} & 1 & 0.0552 & 0.0382 & 0.0791 & 0.0459 & 0.1066 & 0.0528 & 0.0377 & Round reference \\
       & 2 & 0.0443 & 0.0293 & 0.0679 & 0.0369 & 0.0954 & 0.0439 & 0.0295 & Lower scores \\
       & 3 & 0.0562 &\textbf{0.0393}& 0.0783 &\textbf{0.0464}& 0.1075 &\textbf{0.0538}&\textbf{0.0387}& Recovery \\
 \rowcolor{methodblue}      & Ref. &\textbf{0.0563}& 0.0387 &\textbf{0.0800}& 0.0463 &\textbf{0.1087}& 0.0535 & 0.0380 & Validation-selected \\
\hline
\multirow[c]{4}{*}{Toys}   & 1 & 0.0530 & 0.0367 & 0.0749 & 0.0437 & 0.1012 & 0.0504 & 0.0361 & Round reference \\
       & 2 & 0.0521 & 0.0355 & 0.0725 & 0.0421 & 0.0956 & 0.0479 & 0.0344 & Lower scores \\
       & 3 & 0.0490 & 0.0333 & 0.0693 & 0.0399 & 0.0951 & 0.0464 & 0.0326 & Lower scores \\
 \rowcolor{methodblue}      & Ref. &\textbf{0.0592}&\textbf{0.0410}&\textbf{0.0828}&\textbf{0.0487}&\textbf{0.1108}&\textbf{0.0557}&\textbf{0.0401}& Validation-selected \\
\hline
\multirow[c]{4}{*}{Sports} & 1 & 0.0328 & 0.0220 & 0.0496 & 0.0275 & 0.0692 & 0.0324 & 0.0221 & Round reference \\
       & 2 & 0.0332 & 0.0226 & 0.0493 & 0.0278 & 0.0682 & 0.0325 & 0.0225 & Recall--NDCG trade-off \\
       & 3 & 0.0316 & 0.0216 & 0.0472 & 0.0266 & 0.0675 & 0.0318 & 0.0218 & Lower scores \\
 \rowcolor{methodblue}      & Ref. &\textbf{0.0385}&\textbf{0.0265}&\textbf{0.0569}&\textbf{0.0324}&\textbf{0.0779}&\textbf{0.0378}&\textbf{0.0261}& Validation-selected \\
\hline
\end{tabular}
\end{table*}
\FloatBarrier


\Needspace{11\baselineskip}
\section{Conclusion}
\label{sec:conclusion}

RecDataAgent formulates recommendation data augmentation as complete user-conditioned policy construction. A language-model controller analyzes train-visible user and item context, selects users, assigns typed data-construction methods, and allocates their exposure under a fixed budget. Typed execution and an evidence archive then connect the proposed policy to realized records, downstream recommendation outcomes, and subsequent policy revision.

Across Amazon Beauty, Toys, and Sports, the validation-selected complete policies achieve the strongest local performance on the four core Recall/NDCG metrics at @10 and @20. Matched routing controls demonstrate the value of contextual allocation, and the first-round feedback ablation characterizes how validation outcomes influence policy construction. The multi-round analysis connects successive policy decisions with their resulting recommendation performance. Together, these results establish a practical framework for complete-policy construction and revision through validation-guided search.

\FloatBarrier
\section*{AI Use Statement}
Generative AI assisted with manuscript writing and LaTeX formatting. The authors reviewed the manuscript and remain responsible for its contents.
\bibliographystyle{iclr2027_conference}
\bibliography{references}
\appendix
\section{Experimental Details}
\label{app:experiments}

This appendix expands the evaluation setup and intervention evidence in Sections~\ref{sec:exp-setup}--\ref{sec:archive-analysis}. It consolidates the protocol, the first-round feedback comparison, a feature-guided case analysis, and a self-contained planner specification.

\subsection{Protocol and Evidence}

\paragraph{Evaluation protocol.}
The evaluation uses Amazon Beauty, Toys, and Sports sequential recommendation. Each example provides a preceding item sequence and predicts the next item's five characteristic words. The fixed populations contain 22,363 Beauty users, 19,412 Toys users, and 35,598 Sports users. The principal comparisons use Qwen3-4B-Instruct-2507 with one epoch of full-parameter BF16 training, cutoff length 2,048, cosine decay from $10^{-4}$, and effective global batch size 80. DeepSpeed ZeRO-3 is used for training; hardware-dependent micro-batches vary while the effective global batch remains fixed. The five training seeds are $\{40,41,42,43,44\}$ and the reference data seed is 42.
The controlled comparison treats all baseline methods included in Table~\ref{tab:overall} as method-level comparators. Each baseline retains its method-specific construction and recommendation objective, while the shared recommendation setting and one-epoch training horizon keep the downstream task fixed, allowing the experiments to focus on user-conditioned data-policy construction and feedback-guided revision.

For the reported local results, deterministic decoding uses batch size 1, 20 beams, 20 returned sequences, and at most 30 new tokens. All reported local tables and comparisons are restricted to this fixed 20-beam evaluation protocol. Formal results cover six disjoint evaluation shards and are released only after aggregate identity, protocol, and completeness checks pass. Training data provide the base corpus, augmentation inputs, and user features; validation outcomes supply all feedback for proposal revision, candidate ranking, incumbent updates, and final policy selection. Local result tables report metrics after policies and comparison conditions are fixed, while published comparator rows retain their source protocols. Policy-to-run linkage has been checked against the planner prompt, generated policy, materialized corpus, completed training record, and evaluation summary.

\paragraph{Main-table aggregation.}
For each domain, every metric in the locally evaluated rows of Table~\ref{tab:overall}---L2Aug, GenPAS, AsarRec, and \method{}---is the arithmetic mean across the five training seeds $\{40,41,42,43,44\}$. Published comparator rows retain the reporting conventions of their source studies.

\paragraph{Search budget and stopping rule.}
For each domain, policy search runs for three rounds with three candidate policies per round, yielding nine candidates in total. We use a fixed three-round stopping rule: search terminates after completing Round~3. The final policy is selected by validation utility from the completed candidates within this budget. The reference row in Table~\ref{tab:three-rounds} reuses the final selected policy from this search.
Here validation utility is validation Recall@10 from the fixed seed-42 reference run, computed as the user-count-weighted micro-average over complete validation shards. The five-seed arithmetic means in the main table are reporting aggregates for policies fixed before that evaluation, not an additional selection stage. This rule is fixed before candidate generation and is independent of the candidate's declared hypothesis metric; an exact tie is resolved by earlier round and then lower candidate index. Test and Full-evaluation outcomes are not used for candidate ranking or final policy selection.

\paragraph{Round-level selection trace.}
To make the link between search-time selection and the reported round profiles explicit, Table~\ref{tab:round-selection-trace} places the fixed validation utility beside the corresponding Test Recall@10 already reported in Table~\ref{tab:three-rounds}. The validation utility is the seed-42 validation Recall@10 defined above; Test Recall@10 is shown only for audit and does not enter selection. The \textsc{Ref.} row is the final policy selected from the complete nine-candidate archive, not a fourth round. Because Table~\ref{tab:three-rounds} shows one representative policy per round, the reference may come from any of the other six candidates.

\begin{table}[H]
\centering
\caption{Round-level selection trace. Validation utility is seed-42 validation Recall@10; the Test Recall@10 entries are copied from Table~\ref{tab:three-rounds} for audit and do not enter selection.}
\label{tab:round-selection-trace}
\small
\renewcommand{\arraystretch}{1.05}
\setlength{\tabcolsep}{4.5pt}
\begin{tabular}{llrrl}
\hline
Domain & Entry & Validation R@10 & Test R@10 & Selection role \\
\hline
\multirow[c]{4}{*}{Beauty}
 & Round~1 & 0.087326 & 0.0791 & Displayed round \\
 & Round~2 & 0.074418 & 0.0679 & Displayed round \\
 & Round~3 & 0.086756 & 0.0783 & Displayed round \\
 & \textsc{Ref.} & \textbf{0.088000} & \textbf{0.0800} & Final selected \\
\hline
\multirow[c]{4}{*}{Toys}
 & Round~1 & 0.082465 & 0.0749 & Displayed round \\
 & Round~2 & 0.079533 & 0.0725 & Displayed round \\
 & Round~3 & 0.076854 & 0.0693 & Displayed round \\
 & \textsc{Ref.} & \textbf{0.091080} & \textbf{0.0828} & Final selected \\
\hline
\multirow[c]{4}{*}{Sports}
 & Round~1 & 0.054461 & 0.0496 & Displayed round \\
 & Round~2 & 0.054378 & 0.0493 & Displayed round \\
 & Round~3 & 0.052345 & 0.0472 & Displayed round \\
 & \textsc{Ref.} & \textbf{0.062590} & \textbf{0.0569} & Final selected \\
\hline
\end{tabular}
\end{table}

\paragraph{Worked provenance example: Beauty.}
The Beauty archive used for this worked example contains three sequential evaluated policies, namely Round~1 (\texttt{research-1}), Round~2 (\texttt{research-2}), and Round~3 (\texttt{research-3}). Table~\ref{tab:three-rounds} is a compact cross-domain trace, whereas Table~\ref{tab:beauty-candidate-ledger} gives the complete displayed Beauty round-policy ledger; the two tables therefore serve different reporting purposes. The final reference policy is reported separately as \textsc{Ref.} because it is the selected policy used in the main comparison, rather than a fourth round. Validation feedback was used to form later proposals, whereas the ledger reports formal Full-evaluation outcomes under deterministic 20-beam decoding, not validation utility values.

\begin{table}[H]
\centering
\caption{Beauty round-policy ledger for the worked provenance example. Each displayed row is one completed policy; the metrics are formal Full-evaluation results, not validation utility values. The bold row is the final reference policy used in the main comparison.}
\label{tab:beauty-candidate-ledger}
\scriptsize
\setlength{\tabcolsep}{2.8pt}
\begin{tabular}{llrrrrl}
\hline
Round & Candidate & R@10 & N@10 & R@20 & N@20 & Role \\
\hline
1 & \texttt{research-1} & 0.0791 & 0.0459 & 0.1066 & 0.0528 & \\
2 & \texttt{research-2} & 0.0679 & 0.0369 & 0.0954 & 0.0439 & \\
3 & \texttt{research-3} & 0.0783 & 0.0464 & 0.1075 & 0.0538 & \\
 \textsc{Ref.} & \texttt{main-reference} & \textbf{0.0800} & \textbf{0.0463} & \textbf{0.1087} & \textbf{0.0535} & \textbf{Final Ref.} \\
\hline
\end{tabular}
\end{table}

\paragraph{Candidate corpus budget.}
Each candidate retains one complete copy of the base corpus. The principal configurations append 22,363, 19,412, and 35,572 records in Beauty, Toys, and Sports, respectively, as summarized in Table~\ref{tab:data-stats}; the same domain-specific budgets are used in the matched feedback comparison. Every eligible user contributes at most one appended record, and the 26 unsupported Sports users remain base-only. These settings preserve the exact base-record prefix and fixed method applicability across local comparisons.

\paragraph{Evidence organization.}
The archive selects and compares complete corpus policies using validation utility. The appendix organizes the evidence into results for validation-selected policies, contextual published comparators, a matched first-round feedback ablation, and a feature-guided case analysis. Historical design traces are presented separately in Appendix~\ref{app:policy-revision}.

\subsection{Matched First-Round Feedback Ablation}

\paragraph{Comparison setup.}
Within each domain, the two conditions share the same round, planner model and configuration, search starting point, candidate budget, training data, method catalog, exposure constraints, and training and evaluation settings. The feedback condition receives aggregate validation metrics, rankings, incumbent labels, and outcome-derived mechanism statements; the no-feedback condition omits these channels while retaining train-visible features and construction diagnostics. The executor applies the same validation-selection criterion to both conditions.

The pairing has been verified through the policy, prompt, training, and evaluation records. It fixes the experimental conditions while allowing generated policies to respond to feedback. Table~\ref{tab:feedback-raw} reports the paired Round~1 results; Table~\ref{tab:no-feedback} compares the final validation-selected policy with the no-feedback candidate under the same domain-specific augmentation and one-epoch training budgets.

\begin{table}[H]
\centering
\caption{Matched first-round feedback ablation. Within each domain, the paired rows share the round, candidate budget, planner configuration, source data, corpus budget, training, and evaluation protocol; only validation-feedback availability changes.}
\label{tab:feedback-raw}
\small
\renewcommand{\arraystretch}{1.08}
\setlength{\tabcolsep}{3.4pt}
\rowcolors{2}{promptbluebg}{white}
\begin{tabular}{llrrrrrrr}
\toprule
Domain & Planner context & HR@5 & NDCG@5 & R@10 & N@10 & R@20 & N@20 & MRR@20 \\
\midrule
Beauty & Feedback Round 1 & 0.0552 & 0.0382 & 0.0791 & 0.0459 & 0.1066 & 0.0528 & 0.0377 \\
       & No feedback      & 0.0544 & 0.0380 & 0.0770 & 0.0453 & 0.1056 & 0.0525 & 0.0376 \\
Toys   & Feedback Round 1 & 0.0530 & 0.0367 & 0.0749 & 0.0437 & 0.1012 & 0.0504 & 0.0361 \\
       & No feedback      & 0.0537 & 0.0364 & 0.0756 & 0.0435 & 0.1027 & 0.0503 & 0.0355 \\
Sports & Feedback Round 1 & 0.0328 & 0.0220 & 0.0496 & 0.0275 & 0.0692 & 0.0324 & 0.0221 \\
       & No feedback      & 0.0308 & 0.0211 & 0.0456 & 0.0259 & 0.0631 & 0.0303 & 0.0211 \\
\bottomrule
\end{tabular}
\rowcolors{1}{}{}
\end{table}

Under the matched Round~1 conditions, validation feedback improves all seven displayed metrics in Beauty and Sports. In Toys, it improves NDCG@5, NDCG@10, NDCG@20, and MRR@20, while the no-feedback policy is higher on HR@5, Recall@10, and Recall@20. The matched ablation therefore shows gains across all displayed metrics in Beauty and Sports and a Recall--NDCG trade-off in Toys.

\subsection{Feature-Guided Case Study}

\paragraph{Planner-visible feature prompt.}
The planner received the following feature block as part of its Beauty research prompt. Here $L$ is the train-visible history length, a leaf category is the final normalized category component, and each item TID contains five words.

\begingroup\footnotesize
\promptlabel{USER\_FEATURES (name = computation; planner-visible meaning).}
\begin{enumerate}[leftmargin=*,itemsep=1.5pt,topsep=3pt,label=\textcolor{promptblue}{\arabic*.}]
\item \texttt{history\_length} $=L$; available sequential context.
\item \texttt{p\_law\_item\_actual\_entropy} $=\operatorname{actual\_entropy}(\text{item-ID sequence})$; the supplied P-Law statistic, equal to $\ln L$ for this population.
\item \texttt{category\_actual\_entropy} $=\operatorname{actual\_entropy}(\text{ordered leaf-category sequence})$; ordered category complexity.
\item \texttt{category\_unique\_ratio} $=\#\text{unique leaf categories}/L$; category breadth.
\item \texttt{category\_top\_share} $=\max_c \#c/L$; concentration in the most frequent category.
\item \texttt{category\_shannon\_entropy\_norm} $=-\sum_c p_c\ln p_c/\ln L$; normalized category-frequency diversity.
\item \texttt{category\_switch\_rate} $=\#\{c_t\neq c_{t-1}\}/(L-1)$; adjacent category switching.
\item \texttt{tid\_unique\_ratio} $=\#\text{unique TID words}/(5L)$; TID vocabulary breadth.
\item \texttt{tid\_top\_share} $=\max_w \#w/(5L)$; concentration of the most frequent TID word.
\item \texttt{tid\_shannon\_entropy\_norm} $=-\sum_w p_w\ln p_w/\ln(5L)$; normalized TID-word frequency diversity.
\item \texttt{adjacent\_tid\_jaccard\_mean} $=\operatorname{mean}_t |T_{t-1}\cap T_t|/|T_{t-1}\cup T_t|$; local semantic-word continuity.
\item \texttt{early\_late\_tid\_jaccard} $=|T_{\mathrm{early}}\cap T_{\mathrm{late}}|/|T_{\mathrm{early}}\cup T_{\mathrm{late}}|$; persistence between the first and second history halves.
\item \texttt{late\_new\_tid\_ratio} $=|T_{\mathrm{late}}\setminus T_{\mathrm{early}}|/|T_{\mathrm{late}}|$; recently introduced TID words.
\item \texttt{item\_log\_popularity\_mean} $=\operatorname{mean}_i\ln(1+\operatorname{freq}_{\mathrm{train}}(i))$; average item popularity.
\item \texttt{item\_log\_popularity\_std} $=\operatorname{std}_i\ln(1+\operatorname{freq}_{\mathrm{train}}(i))$; within-history popularity range.
\item \texttt{long\_tail\_item\_ratio} $=\#\{i:\operatorname{freq}_{\mathrm{train}}(i)\leq 8\}/L$; share of items in the frozen training-frequency tail.
\item \texttt{popularity\_trend} $=\operatorname{mean}_{\mathrm{late}}\ln(1+\operatorname{freq})-\operatorname{mean}_{\mathrm{early}}\ln(1+\operatorname{freq})$; movement toward more or less popular items.
\end{enumerate}

\promptlabel{CLUSTERING.}
Retain all 17 features in the planner input. Fit the user clusters with the following 14 nonredundant coordinates:
\begin{itemize}[leftmargin=*,itemsep=0pt,topsep=2pt,parsep=0pt]
\item \texttt{history\_length}, \texttt{category\_actual\_entropy}, \texttt{category\_unique\_ratio}, \texttt{category\_top\_share}.
\item \texttt{category\_shannon\_entropy\_norm}, \texttt{category\_switch\_rate}, \texttt{tid\_unique\_ratio}, \texttt{tid\_top\_share}.
\item \texttt{adjacent\_tid\_jaccard\_mean}, \texttt{early\_late\_tid\_jaccard}.
\item \texttt{item\_log\_popularity\_mean}, \texttt{item\_log\_popularity\_std}.
\item \texttt{long\_tail\_item\_ratio}, \texttt{popularity\_trend}.
\end{itemize}
Apply median/IQR RobustScaler normalization. Compare KMeans $K=2,\ldots,8$ with \texttt{random\_state=42} and \texttt{n\_init=20}; among candidates whose smallest cluster contains at least 2\% of users, select the highest sampled silhouette.
\endgroup

\paragraph{Meaning of \texttt{cluster\_id}.}
The resulting \texttt{cluster\_id} is the integer label produced by this frozen Beauty-specific KMeans fit. Figure~\ref{fig:beauty-clusters} shows why $K=2$ was selected and how the two groups differ. Cluster~0 contains 18,175 users and is higher on category breadth and switching. Cluster~1 contains 4,188 users and is higher on category concentration and TID-word overlap. Thus, \texttt{cluster\_id == 1} is a concrete predicate selecting the second feature group; it is not a user identifier or a rank.

\begin{figure*}[t]
\centering
\includegraphics[width=\textwidth]{figures/beauty_cluster_analysis.pdf}
\caption{Beauty feature geometry used for policy routing. \textbf{A:} Sampled silhouette scores across candidate values of $K$ identify the selected cluster count. \textbf{B:} PCA projects the planner features with cluster assignments and projected medians. \textbf{C:} Cluster medians are shown relative to the full-population median.}
\label{fig:beauty-clusters}
\end{figure*}

\Needspace{18\baselineskip}
\paragraph{Representative augmentation example.}
Consider a train-visible Beauty history of length four assigned to Cluster~1. Its feature profile belongs to the group with stronger category concentration and adjacent TID-word overlap, for which the selected policy supplies a GenPAS prefix view. The following excerpt shows how that history is presented to the recommender after augmentation.

\begingroup\footnotesize
\begin{promptblock}[samepage=true]
Item text ID: [nail, art, pearl, stone, small]
Title: Islandoffer Color Pearl Nail Art Stone Small Wheel
Rhinestones Beads.

<genpas_augmented_sequence>
Item ID: 6066; Item text ID: [nail, art, pearl, stone, small]
Item ID: 4171; Item text ID: [nail-art, 3d, colorful,
slices, diy]
Item ID: 1270; Item text ID: [nail-art, rhinestones,
white-pearl, round, 1.5mm]
\end{promptblock}
\endgroup

The original recommendation instruction and target remain unchanged; the augmentation contributes an additional structured view of the observed history. This example illustrates the policy-level transformation without introducing a separate user-level label or evaluation target.

\subsection{Planner Prompt and Feature-Guided Policy}

\paragraph{Prompt organization.}
The formal planner input combines the task, train-visible data and features, clustering analysis, method definitions, applicability constraints, released validation evidence, and the JSON response contract. Table~\ref{tab:prompt-records} summarizes these components. The accompanying \emph{RecDataAgent Planner Protocol and Prompt Specification} provides an author-edited English specification of the formal interface, including the feedback-conditioned and no-feedback settings.

\begin{table}[H]
\centering
\caption{Planner-input specification for feature-guided policy construction. Task context, user features and groups, method execution context, and validation evidence define the input channels; validation evidence supplies the outcome signal for policy revision.}
\label{tab:prompt-records}
\small
\renewcommand{\arraystretch}{1.14}
\setlength{\tabcolsep}{4pt}
\rowcolors{2}{promptbluebg}{white}
\begin{tabular}{p{0.29\columnwidth}p{0.61\columnwidth}}
\toprule
Component & Content and role \\
\midrule
Task and data context & Recommendation task, train-visible histories, item information, and the unchanged base corpus. \\
Features and user groups & Feature definitions, distribution summaries, clustering, and executable route predicates. \\
Method and execution context & Baseline methods, applicability, exposure constraints, and the response schema. \\
Validation feedback & Completed validation outcomes, rankings, and incumbent labels for policy revision. \\
\bottomrule
\end{tabular}
\rowcolors{1}{}{}
\end{table}

\noindent The no-feedback condition retains the task, feature, and execution components while withholding outcome-derived planner context. The executor applies the same fixed validation-selection criterion in both conditions.

\paragraph{Language and presentation.}
The production planner prompt was issued in Chinese. The feature definitions and operative instructions are organized here as an English protocol specification, with the Beauty policy providing a concrete case study.

\paragraph{Recommendation instruction.}
\begingroup\footnotesize
\begin{promptblock}
Based on the user's historical product interaction sequence,
predict the next product's characteristic words.
Each product is represented by exactly 5 characteristic words
enclosed in square brackets []. The historical sequence shows
the user's interaction pattern.
\end{promptblock}
\endgroup

\paragraph{Operative prompt rendering.}
\begingroup\footnotesize
\noindent\textbf{Role and task.}
Act as a recommendation-data research agent for Beauty sequential recommendation with Qwen3-4B-Instruct-2507. Propose exactly three substantively different augmentation hypotheses, rather than three variants of one length threshold. Separate observed statistics, mechanism hypotheses, and candidate evaluation objectives.

\noindent\textbf{Data and input.}
The base corpus contains 34,464 JSON records, each with exactly three string fields: \texttt{instruction}, \texttt{input}, and \texttt{output}. Its first 22,363 records are recommendation examples and its remaining 12,101 records are item-description-to-five-word-TID examples. The augmentation source contains 22,363 users with train-visible histories of length 3--18. Policy generation is restricted to \texttt{user\_idx}, \texttt{sequence\_item\_ids}, \texttt{sequence\_tids}, derived train-visible features, and the train-visible catalog.

\noindent\textbf{Base corpus and appended views.}
Each candidate dataset retains one unchanged copy of the complete base corpus in its original order. An augmentation creates a new record by copying the matched base record's \texttt{instruction} and \texttt{output} and appending the selected method view to its \texttt{input}. Each selected user contributes one appended row, giving coverage histogram \texttt{\{"1": N\}} for $N$ unique users.

\noindent\textbf{Available methods and routes.}
Use the listed methods with their actual signatures. The available method implementations are \texttt{asarrec\_dual\_view}, \texttt{genpas\_prefix\_view}, and \texttt{l2aug\_sequence\_view}. A feature route contains one to eight conjunctive numeric predicates using \texttt{<}, \texttt{<=}, \texttt{>}, \texttt{>=}, \texttt{==}, or inclusive \texttt{between}; \texttt{cluster\_id} may only be compared for equality to an existing integer. Within a candidate, routes must be mutually exclusive. The selected streams are combined by the declared route assignment, and any uncovered selected user must go to an explicit remainder stream.

\noindent\textbf{Coverage conditions.}
Execution checks method applicability for every selected user, preserves unique-user exposure, and routes the complete remainder explicitly. A missing predicate or method triggers a minimal capability request and dry-run before materialization.

\noindent\textbf{Evidence checks.}
Exact joint-predicate counts, interface signatures, and applicability are verified through \texttt{code\_evidence}. Each dry-run reports covered users, remainder users, cross-route overlap, and method applicability.

\noindent\textbf{Feedback and selection.}
Use aggregate validation outcomes to revise proposals. The experiment-level validation objective is fixed before candidate generation and governs ranking and final selection. Candidate-specific hypotheses describe comparisons within that protocol. Test outcomes are reserved for reporting after the policies and comparisons are fixed.

\Needspace{25\baselineskip}
\noindent\textbf{Response protocol.}
Return a complete JSON object whose root contains exactly \texttt{action}, \texttt{reason}, \texttt{experiments}, and \texttt{requests}. The action is either \texttt{execute} or \texttt{needs\_capability}. For \texttt{execute}, return exactly three experiments and an empty request list; for \texttt{needs\_capability}, return an empty experiment list and at least one concrete request.

\begin{promptblock}[samepage=true]
{
  "action": ..., "reason": ...,
  "experiments": [{
    "name": ..., "total_exposure": ...,
    "operators": [{"name": ..., "params": ...}],
    "rationale": ..., "risks_and_checks": ...,
    "coverage": {
      "unique_users": N,
      "exposure_histogram": {"1": N}
    },
    "comparison": {
      "reference": ..., "primary_metric": ...,
      "hypothesis": ..., "changed": [...], "fixed": [...],
      "interpretation": ...
    }
  }],
  "requests": [{
    "kind": ..., "target": ...,
    "details": ..., "why_needed": ...
  }]
}
\end{promptblock}
Each request \texttt{kind} must be either \texttt{implement\_function} or \texttt{code\_evidence}; each method entry contains exactly \texttt{name} and \texttt{params}. The primary metric must be one of Recall@10, NDCG@10, Recall@20, or NDCG@20.
The \texttt{primary\_metric} field is a candidate-level hypothesis annotation rather than a selection control: it may name the metric emphasized by that candidate's proposed comparison, but it cannot override the fixed experiment-level validation utility, which is validation Recall@10 as specified above.
\endgroup

\paragraph{Feature-guided policy definition.}
\begingroup\footnotesize
\begin{promptblock}
genpas_prefix_view(
  source_id="beauty-all-train-visible-v1",
  policy_id="genpas_full_available_seed42_v1",
  output_stream="cgeom_genpas",
  applicability="require_all")

asarrec_dual_view(
  source_id="beauty-all-train-visible-v1",
  policy_id="asarrec_full_available_seed42_v1",
  output_stream="cgeom_asarrec",
  applicability="require_all")

feature_route_select(
  routes=[{
    stream: "cgeom_genpas",
    query: {all: [{
      feature: "cluster_id", op: "==", value: 1
    }]},
    applicability: "require_all"
  }],
  remainder_stream="cgeom_asarrec",
  merge="concatenate")

emit_standard_artifacts()
\end{promptblock}
\endgroup

\paragraph{Alternative hypotheses.}
The response also proposed routes based on zero early-versus-late TID-word overlap and on joint long-tail exposure with decreasing popularity. The selected case above records the validation-selected cluster-geometry policy used for the reported Beauty result.

\paragraph{Research linkage.}
The prompt, feature representation, clustering analysis, and representative augmentation example define the intervention studied in the Beauty experiment. The policy specification maps train-visible feature geometry to the available augmentation views while preserving the original recommendation supervision. The resulting corpus is evaluated under the validation-guided selection protocol in Section~\ref{sec:exp-setup}, and the selected policy's Beauty performance is given in Table~\ref{tab:overall}. Together, these elements connect the planner input, the policy decision, and the corpus-level experiment.

\Needspace{8\baselineskip}
\section{Feedback-Driven Policy Revision}
\label{app:policy-revision}

This section documents how feedback informs both the analytical basis
used to route users and the augmentation configuration applied after
routing. The historical design records link planner inputs, method
plans, realized coverage, and recorded aggregate outcomes. They complement
the matched feedback ablation in Table~\ref{tab:feedback-raw} and are
presented separately from the formal rounds in Table~\ref{tab:three-rounds}.

\paragraph{What counts as a revision.}
We treat route predicates and feature partitions as the \emph{analytical
basis}, and methods, exposure allocation, merge rule, and repetition as the
\emph{augmentation configuration}. This separates whom to route from which
view to supply and at what dose.

\begin{table}[H]
\centering
\caption{Historical Beauty policy revisions and recorded outcomes. Each trace lists its routing basis, method exposure, and Recall@10; the matched feedback ablation is reported separately in Table~\ref{tab:feedback-raw}.}
\label{tab:policy-revision-trace}
\small
\renewcommand{\arraystretch}{1.08}
\setlength{\tabcolsep}{3pt}
\begin{tabular}{@{}>{\raggedright\arraybackslash}p{0.15\columnwidth}
                >{\raggedright\arraybackslash}p{0.34\columnwidth}
                >{\raggedright\arraybackslash}p{0.36\columnwidth}
                r@{}}
\toprule
Trace & Routing basis & Method / dose & R@10 \\
\midrule
research-1 &
GenPAS: $L=3$--6; AsarRec: $L=7$--18 &
GenPAS 15,874; AsarRec 6,489; total 22,363 &
0.0771 \\
research-2 &
AsarRec: $L=3$--5, $10$--18; GenPAS: $L=6$--9 &
AsarRec 17,610; GenPAS 4,753; total 22,363 &
0.0626 \\
research-3 &
GenPAS: $L=3$--6, $13$--18; AsarRec: $L=7$--12 &
GenPAS 18,063; AsarRec 4,300; total 22,363 &
0.0736 \\
feedback candidate &
No length route; stable round-robin &
AsarRec 11,875; GenPAS 3,125; total 15,000; 3,125 repeated &
0.0752 \\
\bottomrule
\end{tabular}
\end{table}

\paragraph{Direct evidence from the trace.}
The analytical basis is not fixed: the first trace uses two length segments,
the second uses three with GenPAS in the middle, and the third moves
$L=13$--18 back to GenPAS while retaining AsarRec for $L=7$--12. The recorded
predicates document successive changes in contextual assignment.

The augmentation configuration also changes. Across the three research
traces, the method pair remains GenPAS plus AsarRec, but the exposure
allocation changes from 15,874/6,489 to 4,753/17,610 and then to
18,063/4,300. The feedback candidate changes the budget to 15,000 and uses
stable round-robin mixing with 3,125 repeated exposures. Thus, both routing
and method/dose configuration are revised.

\paragraph{Feedback enters the next design.}
The feedback-conditioned record compares historical single-method results,
notes the all-user AsarRec baseline (Recall@10 $=0.0781$) above GenPAS ($=0.0733$)
and L2Aug ($=0.0677$), and proposes an AsarRec-heavy mixture with a smaller
GenPAS component. That candidate materialized with 15,000 appended rows, one
unchanged base-corpus copy, and zero target or view truncation. This record shows
how archived outcomes inform the next policy's method and exposure choices.

\paragraph{Relation to the formal experiments.}
The verified records organize this evidence by its experimental role.
These design examples describe changes in routing and exposure.
Table~\ref{tab:feedback-raw} provides the matched feedback comparison, and
Table~\ref{tab:three-rounds} reports the formal validation-guided sequence.

\end{document}
